\section{Binding}
\label{sec:binding}

Motion under the candidate is a one-body statement. Matter is bound, and a bound composite must hold together, move
isotropically under $c^*$, and have inertia equal to its energy. Under flat links separation costs nothing, and a level
bound by contact alone sits inside the free two-hole continuum and decays (\code{E-SPN-0145}). The string must be
supplied, and each attempt below narrows what it can be.

\subsection{A cost on separation}

The string closes only on charge 0 mod 3, so two holes cannot carry one. A hole on the love sea is exactly a love on
the empty mesh, so the test is the love-fear meson on the empty mesh.
\begin{itemize}
\item \textbf{A phase string} (angle 0.0936 per unit of Steiner length, exact in $\mathbb{Z}[\omega][1/42]$) forms a
composite at energy 1.5658 against 1.5705 predicted, isotropic to $9.5 \times 10^{-9}$, $R = 1.027$, moving at 0.17,
0.32 and $0.41c^*$ at $K$ = 0.4, 0.8 and 1.2. It is a resonance: at length $2m/\sigma$ the string has paid $2m$, and a
member drops into 11 flat bands degenerate with its partner, where no Schwinger suppression applies
\cite{kurzynski2008} (\code{E-SPN-0146}).
\item \textbf{A mass string}, the mixer angle growing with separation, closes that channel: the mixer acts only on a
dock's uniform mode, the flat bands stay pinned at $\pi$, and every flipped channel tops out at $\pi - 2m_0$. The level
holds exactly, fidelity $1 - 1.7 \times 10^{-11}$ over 256 beats with $1.8 \times 10^{-18}$ reaching the edge, and moves
isotropically. But it needs $m_0 > \arctan(1/3)$, reads the separation nonlocally, and gives $R = 1.568$
(\code{E-SPN-0147}).
\end{itemize}

\begin{theorem}[The band sum]
For the stream times any fixed dock matrix, the trace has no constant Fourier term because no root is zero, so the
band phases sum to zero over the zone. At the swap point the moving pair averages to zero, so if the pair and the even
flat bands are kept, the 11 odd flat bands stay at $\pi$. No dock-local unitary, in any ring, moves them alone
(\code{E-SPN-0148}).
\end{theorem}

The motive was also wrong: the level carries only 0.16\% per member in the flat bands, and $R$'s excess lives in the
moving pair. On a line, where Gauss's law fixes the flux and a local string equals the instantaneous one, $R \to 1$ with
the lattice mass $\tan(m)/m$ the only excess, falling as $n^{-2}$ (measured slope $-2.078$), and binding moves $R$ by at
most $1.4 \times 10^{-3}$ (\code{E-SPN-0149}). For an instantaneous string in $d$ dimensions the inertia is
$2m + (1 + 2/d)T$ against a rest energy of $2m + 3T$, so a cost-only string misses $(2 - 2/d)T$, which is the
transverse inertia of a real string. The string needs dynamics of its own.

\subsection{A gauge field on the links}

\begin{itemize}
\item \textbf{$\mathbb{Z}_3$ flux.} A local flux register is exact, Gauss exact and vacuum inert, but it records each
member's path, so a member tied to a charge walks the 24-regular tree, the universal cover of the $D_4$ root graph, and
sits above the Kesten floor $\pi/2 - \arcsin(\rho \cos m_0)$, $\rho = 2\sqrt{23}/24$, at least 1.159 for every $m_0$
\cite{kesten1959}. Measured edges match the floor to $10^{-3}$ (\code{E-SPN-0150}).
\end{itemize}

\begin{theorem}[A]
A magnetic plaquette piece is diagonal in link values and the hop is controlled by link values, so the two commute for
every gauge group. The plaquette moves a member only through the electric piece (\code{E-SPN-0151},
\code{E-SPN-0152}).
\end{theorem}

\begin{theorem}[B]
In the flat sector a tied member walks $D_4$ exactly with its edge at $m_0$, but there every link's flux is uniform, so
for any finite group the string tension is exactly 0 (\code{E-SPN-0151}, \code{E-SPN-0153}).
\end{theorem}

Light and confined therefore needs a continuous confinement transition. $\mathbb{Z}_3$ in 4d has a first-order one
\cite{creutz1979}, and in the $4+1$ bulk even $SU(2)$ does \cite{creutz1979su2}, so the string must live on the husk.
The group the roots already form is $2T$: a fixed linear map sends the 24 roots onto the Hurwitz units, all 7
characters lie in $\mathbb{Z}[\omega]$, and a member carries the spinor 2. But $2T$ freezes while the string still has
0.82 of its strong-coupling energy per link, and $2O$ at 0.47 to 0.65 (\code{E-SPN-0152}), as finite subgroups of
$SU(2)$ do \cite{petcher1980, bhanot1981}.

\subsection{2I on the husk}

The husk lattice is cubic with 6 axis and 12 face-diagonal links per dock and 20 triangles, isotropic when a beat is
$\sqrt{3/5}$ of an axis link. A deterministic one-loop crossing of $SU(2)$'s Gaussian free energy with the frozen branch
puts $2I$'s hypercubic freezing at $\beta = 5.40$ against Monte Carlo's $5.70(20)$ \cite{hartung2022}. On the husk
$2I$ freezes at $\beta = 2.82$, matched hypercubic $\beta = 9.77$, where the string tension is $1.6 \times 10^{-18}$ of
its strong-coupling value (\code{E-SPN-0153}).

\begin{result}
A bounded 120-value register reaches $SU(2)$'s scaling region on the husk and loosens the string continuously. No
continuous group is needed. The ring grows to $\mathbb{Z}[\omega, \varphi][1/210]$.
\end{result}

What fails is keeping. No ordered $2I$ register state is kept by the beat: the flat register keeps 0.813 per link per
beat and the empty one 0.630 per face. A beat is expected to keep only locally maximally mixed states. What is proved
is prethermal keeping \cite{abanin2017}: at small beat angle $\delta$ a bounded register holds a state for time
$\exp(c/\delta)$. Whether that holds the ordered $2I$ vacuum long enough is the one open question left in binding, and
it is being run.
